% ---------------------------------------------------------------
% Präambel
% ---------------------------------------------------------------
\usepackage[T1]{fontenc}
\usepackage[utf8]{inputenc}       % ab LaTeX 2018 Standard, schadet nicht
\usepackage[ngerman]{babel}
\usepackage{lmodern}
\usepackage{microtype}
\usepackage{csquotes}

\usepackage[a4paper, left=30mm, right=25mm, top=25mm, bottom=30mm]{geometry}

% Mathematik, Einheiten
\usepackage{amsmath, amssymb}
\usepackage[locale=DE, per-mode=symbol]{siunitx}

% Tabellen, Abbildungen
\usepackage{graphicx}
\graphicspath{{img/}{../matlab/reports/img/}{../docs/img/}}
\usepackage{booktabs}
\usepackage{tabularx}
\usepackage{float}
\usepackage[font=small, labelfont=bf]{caption}
\usepackage{subcaption}

% Farben, Quelltext
\usepackage[dvipsnames]{xcolor}
\usepackage{listings}
\lstset{
  basicstyle=\ttfamily\small,
  breaklines=true,
  frame=single,
  numbers=left,
  numberstyle=\tiny\color{gray},
  keywordstyle=\color{NavyBlue}\bfseries,
  commentstyle=\color{ForestGreen},
  stringstyle=\color{BrickRed},
  showstringspaces=false,
  literate={ä}{{\"a}}1 {ö}{{\"o}}1 {ü}{{\"u}}1 {Ä}{{\"A}}1 {Ö}{{\"O}}1 {Ü}{{\"U}}1 {ß}{{\ss}}1,
}
\lstdefinestyle{matlab}{language=Matlab}
\lstdefinestyle{c}{language=C}

% Abkürzungen (printonlyused erst setzen, wenn im Text \ac{...} verwendet wird,
% sonst bricht die leere Liste die Übersetzung ab)
\usepackage{acronym}

% Literatur
\usepackage[backend=biber, style=numeric, sorting=none]{biblatex}
\addbibresource{literatur.bib}

% Verweise (hyperref zuletzt, cleveref danach)
\usepackage[hidelinks, pdfusetitle]{hyperref}
\usepackage[ngerman, noabbrev]{cleveref}

% ---------------------------------------------------------------
% Eigene Befehle
% ---------------------------------------------------------------
% Platzhalter / noch zu schreiben (im PDF rot sichtbar)
\newcommand{\offen}[1]{\textcolor{red}{[#1]}}
% Datei-, Skript- und Blocknamen
\newcommand{\datei}[1]{\texttt{#1}}
\newcommand{\skript}[1]{\texttt{#1}}
\newcommand{\block}[1]{\textsf{#1}}
% Verweis auf den Hardwareteil des Laborpartners
\newcommand{\hwteil}[1]{(siehe Hardwareteil, #1)}
